\chapter{Options~: arbitrages entre plateformes et sur volatilité à grande vitesse}\label{ch:hf:options-cross-venue-and-volatility-arbitrage-at-speed}

La veille de la date de détachement du dividende d'une action, des milliers de calls profondément dans la monnaie devraient être exercés, et certains ne le sont pas~; les détenteurs qui oublient paient les
ceux qui y pensent, et le trade qui le perçoit figure dans la littérature académique depuis près de vingt ans. Pool, Stoll et Whaley ont constaté
que plus de la moitié des positions longues qui auraient dû être exercées ne l'ont pas été, sur dix ans jusqu'en 2006, à un coût pour les détenteurs supérieur à 491
millions de dollars. Dans le modèle de ce chapitre, le dividend play commence à rapporter lorsque 1.7~\% des détenteurs oublient~; à 10~\%, il rapporte 1.05~\$ par contrat public ouvert,
à 50~\%, 10~\$.

\section{Parité à grande vitesse~: conversions et reversals}

La parité call-put (livre~1, chapitre~25) relie un call, un put et le forward. Lorsque les cotations la violent au-delà des coûts, trois jambes verrouillent la différence.

\begin{definition}[Arbitrage de conversion, arbitrage de reversal]\label{def:hf:options-cross-venue-and-volatility-arbitrage-at-speed:conversion}
Un \emph{arbitrage de conversion}\index{arbitrage de conversion} achète l'action, achète un put et vend un call de même prix d'exercice et même échéance, en verrouillant le
prix d'exercice à l'échéance, lorsque le call est cher par rapport au put et à l'action. Un \emph{arbitrage de reversal}\index{arbitrage de reversal} fait l'inverse~:
vend à découvert l'action, vend le put et achète le call, lorsque le call est bon marché~; il paie les frais d'emprunt de l'action et les éventuels dividendes.
\end{definition}

Lorsque plusieurs plateformes cotent la même série, le meilleur cours acheteur et le meilleur cours vendeur de chaque option, pris sur l'ensemble des plateformes, peuvent violer la parité alors qu'aucune plateforme ne le fait isolément.
Dans le scan de ce chapitre (\cref{lst:hf:options-cross-venue-and-volatility-arbitrage-at-speed:scan}), neuf prix d'exercice d'une chaîne à trois mois sont cotés avec une largeur de dix
cents sur chaque plateforme, autour des fair values issues de Black-76, le mid de chaque plateforme étant décalé de son propre bruit de trois cents~; l'action se tient à un cent de part et d'autre de
100~; chaque jambe coûte un cent. Une violation apparaît dans 0.7~\% des instantanés avec une seule plateforme, 6.6~\% avec trois et 26~\% avec six, valant un à deux cents par
action. Puis les cotations se rafraîchissent une fois, en conservant une part de leur bruit~: lorsque la moitié persiste, seules 6--8~\% des violations subsistent~; lorsque 90~\%
persiste, 43~\%. L'arbitrage existe en proportion du nombre de plateformes et survit d'autant mieux que celles-ci se mettent à jour lentement.

\section{Box spreads et les taux qu'ils impliquent}

Une box (livre~5, chapitre~1) combine un bull call spread et un bear put spread de mêmes prix d'exercice~: elle verse $K_2-K_1$ à l'échéance, quoi
qu'il advienne, son prix est donc celui d'une obligation zéro-coupon. Une box entre 95 et 105, à échéance dans trois mois et valant 9.90~\$, prête au taux de 4.02~\% capitalisé en continu~; si
le taux propre de l'arbitragiste est de 4~\%, l'acheter à 9.89~\$ rapporte un cent. Van Binsbergen, Diamond et Grotteria l'ont exploité~: en déduisant les taux sans risque des
prix d'options, sans modèle de risque, ils ont trouvé un rendement d'opportunité (convenience yield) des bons du Trésor d'environ 40 points de base, plus élevé pour les échéances inférieures à trois mois et
quatre fois plus élevé pendant la crise financière. Le taux implicite de la box est un taux de marché pour les prêts garantis par du collatéral via la chambre de compensation, et négocier
la box contre son propre financement n'est un arbitrage que dans la limite des coûts des quatre jambes et du capital que mobilisent les positions.

\begin{definition}[Jelly roll]\label{def:hf:options-cross-venue-and-volatility-arbitrage-at-speed:jelly}
Un \emph{jelly roll}\index{jelly roll} est un forward synthétique long (achat de call, vente de put) à une échéance contre un forward synthétique court à une autre,
au même prix d'exercice~; son prix est la différence de deux forwards et implique le coût de portage (financement diminué des dividendes et des frais d'emprunt) entre les deux
échéances.
\end{definition}

Dans la chaîne de ce chapitre, un \omterm{def:hf:options-cross-venue-and-volatility-arbitrage-at-speed:jelly}{jelly roll} entre les échéances à trois et six mois implique un coût de portage de 4.0~\%, soit le taux du modèle, puisqu'aucun dividende ne tombe
entre elles. Un dividende annoncé pour cette période, ou une modification des frais d'emprunt de l'action, fait bouger le \omterm{def:hf:options-cross-venue-and-volatility-arbitrage-at-speed:jelly}{jelly roll} avant toute autre chose.

\section{Exercice anticipé, dividendes et dividend play}

Les calls américains sur des actions versant un dividende doivent être exercés la veille de la date de détachement lorsque le dividende dépasse la valeur temps abandonnée
(livre~5, chapitre~6). Pour un call de prix d'exercice 40 sur une action à 60~\$ avec un mois restant, un dividende supérieur à 13 cents rend l'exercice optimal. Les détenteurs qui
n'exercent pas laissent leurs contreparties vendeuses non assignées~: celles-ci conservent un call qui a perdu son dividende. L'assignation se répartit au prorata de toutes les positions
vendeuses ouvertes, donc les teneurs de marché qui se cèdent entre eux, juste avant la date de détachement, de grandes positions opposées puis exercent toutes leurs positions longues, captent l'essentiel des
assignations que les détenteurs distraits n'ont pas provoquées (\texttt{firm.optmm.best\_play} du livre~5). Hao, Kalay et Mayhew ont montré que ce schéma
gonflait le volume déclaré le dernier jour cum-dividende, facilité par les limites que les bourses imposent aux frais de transaction répercutés aux traders.

\begin{omfigure}
\begin{tikzpicture}
\begin{axis}[width=11cm,height=5.4cm,axis y line*=right,axis x line=none,xmode=log,log basis x=10,xmin=0.9,xmax=55,ymin=0,ymax=100,
  ylabel={\small valeur captée (\%)},tick label style={font=\footnotesize},table/col sep=comma]
\addplot[omDef,thick,dashed,mark=square*,mark size=1.3pt] table[x=fail,y=captured]{figdata/hft/20-options-cross-venue-and-volatility-arbitrage-at-speed/dividend.csv};
\label{plot:hf:divcap}
\end{axis}
\begin{axis}[width=11cm,height=5.4cm,axis y line*=left,xmode=log,log basis x=10,xmin=0.9,xmax=55,xtick={1,2,5,10,20,50},
  xticklabels={1,2,5,10,20,50},xlabel={\small détenteurs qui n'exercent pas (\%)},ylabel={\small gain net en \$ par contrat public},ymin=0,ymax=11,
  tick label style={font=\footnotesize},grid=major,grid style={gray!25},table/col sep=comma,
  legend style={font=\scriptsize,at={(0.03,0.97)},anchor=north west}]
\addplot[omThm,thick,mark=*,mark size=1.3pt] table[x=fail,y=per_public]{figdata/hft/20-options-cross-venue-and-volatility-arbitrage-at-speed/dividend.csv};
\addlegendentry{profit net (gauche)}
\addlegendimage{/pgfplots/refstyle=plot:hf:divcap}\addlegendentry{part captée (droite)}
\end{axis}
\end{tikzpicture}

\omcaption{Le dividend play à sa taille optimale~: profit net des deux teneurs de marché par contrat de positions ouvertes publiques (10\,000 contrats), et
la part de la valeur laissée par les détenteurs défaillants que ces teneurs captent, face à la part des détenteurs qui n'exercent pas (échelle logarithmique)~; 30 cents
par action gagnés sur chaque call vendeur non assigné, frais de transaction de 0.2 cent et frais d'exercice de 0.1 cent par action. Données~: \texttt{hf\_optarb.dividend}.}\label{fig:hf:options-cross-venue-and-volatility-arbitrage-at-speed:dividend}
\end{omfigure}

Le profit du dividend play croît avec la part qui fait défaut (\cref{fig:hf:options-cross-venue-and-volatility-arbitrage-at-speed:dividend}), et il en va de même pour la
part de la valeur des défaillances que captent les acteurs~: 42~\% à 5~\% de défaillance, 82~\% à la moitié. Sous le seuil de rentabilité, $f<c/(2g)$ avec $c$ les frais par contrat
échangé et $g$ le gain par call non assigné, le dividend play ne rapporte à aucune taille~: 1.7~\% avec les frais du chapitre.

\section{La même volatilité sur plusieurs cotations}

Les options d'un indice se négocient sur l'indice lui-même (réglées en espèces, européennes) et sur le fonds indiciel coté (ETF) qui le réplique (réglées physiquement, américaines, avec
dividendes)~; les options d'une action se négocient sur plusieurs bourses, identiques et fongibles. La première paire devrait afficher la même volatilité après ajustement
pour le style d'exercice et les dividendes~; la seconde, les mêmes prix. Un teneur de marché qui cote les deux les couvre l'une par l'autre, et le spread entre leurs
volatilités implicites, nets des ajustements, est un trade de relative value qui ne dure que quelques secondes lorsque les cotations de l'une retardent sur celles de l'autre
(le lead--lag du chapitre~8, appliqué à la volatilité).

\section{Fiches de stratégie}

\begin{strategyfile}{Arbitrage de conversion et de reversal}\label{strat:hf:options-cross-venue-and-volatility-arbitrage-at-speed:conversion}
\sfield{Qui paie, et pourquoi} Les plateformes dont les cotations retardent sur celles des autres, et les traders qui traversent une fourchette large sur une seule plateforme.
\sfield{Instruments et plateformes} Une action, ses calls et ses puts sur chaque bourse d'options.
\sfield{Signal} Violations de la parité entre les meilleures cotations des plateformes, au-delà des coûts des trois jambes, des frais d'emprunt et des dividendes.
\sfield{Dimensionnement et exécution} Les trois jambes en même temps, sur les options des plateformes qui cotent la violation~; conserver jusqu'à l'échéance ou procéder au débouclage lorsque la parité revient.
\sfield{Coûts} Les frais des trois jambes, les frais d'emprunt pour les reversals, le risque d'assignation anticipée sur les options américaines.
\sfield{Comment elle s'éteint} Vitesse~: dans le modèle, avec des cotations conservant la moitié de leur bruit, seules 6--8~\% des violations survivent à un rafraîchissement.
\sfield{Horizon, capacité, infrastructure} De la microseconde à l'échéance~; le flux de chaque plateforme.
\sfield{Backtester honnêtement} Cotations telles que reçues, rafraîchissement qui suit, frais d'emprunt et dividendes tels qu'ils étaient connus à l'époque.
\sfield{Sources} Livre~1, chapitre~25~; ce chapitre.
\end{strategyfile}

\begin{strategyfile}{Arbitrage de financement par box spread}\label{strat:hf:options-cross-venue-and-volatility-arbitrage-at-speed:box}
\sfield{Qui paie, et pourquoi} Les emprunteurs et les prêteurs qui utilisent des box pour se financer, et les détenteurs de bons du Trésor qui paient pour cette commodité.
\sfield{Instruments et plateformes} Options sur indice de style européen (sans exercice anticipé), quatre jambes.
\sfield{Signal} Le taux implicite de la box comparé au financement de la firme et aux taux du Trésor.
\sfield{Dimensionnement et exécution} Prêter ou emprunter par la box lorsque son taux dépasse celui du financement d'un écart supérieur aux coûts des quatre jambes.
\sfield{Coûts} Frais et fourchettes des quatre jambes, marge, bilan.
\sfield{Comment elle s'éteint} C'est un trade de financement à marge étroite~; le rendement d'opportunité est le prix que le marché fait payer pour la sécurité des bons du Trésor, et non un
repas gratuit.
\sfield{Horizon, capacité, infrastructure} De quelques semaines à plusieurs années~; capacité fixée par la profondeur du marché des box.
\sfield{Backtester honnêtement} Le taux de financement (funding rate) effectif de la firme et les charges de capital.
\sfield{Sources} van Binsbergen, Diamond et Grotteria (2022).
\end{strategyfile}

\begin{strategyfile}{Dividend play sur calls non exercés}\label{strat:hf:options-cross-venue-and-volatility-arbitrage-at-speed:dividend}
\sfield{Qui paie, et pourquoi} Les détenteurs de calls qui n'exercent pas avant la date de détachement.
\sfield{Instruments et plateformes} Calls profondément dans la monnaie sur des actions versant un dividende, le dernier jour de cotation cum-dividende.
\sfield{Signal} Calls dont l'exercice optimal est certain (dividende supérieur à la valeur temps), et positions ouvertes.
\sfield{Dimensionnement et exécution} Négocier de grandes positions compensatoires avec un autre intervenant, exercer toutes les positions longues~; dimensionner selon la part attendue d'échecs.
\sfield{Coûts} Frais de transaction et d'exercice (que les plafonds des bourses limitent)~; le risque que tout le monde exerce.
\sfield{Comment elle s'éteint} Les détenteurs apprennent à exercer, changements de frais, règles encadrant ce type de transactions.
\sfield{Horizon, capacité, infrastructure} Un jour par date de détachement~; capacité fixée par les positions ouvertes non exercées.
\sfield{Backtester honnêtement} Attribution selon le processus réel de la chambre de compensation~; frais tels que plafonnés à l'époque.
\sfield{Sources} Pool, Stoll et Whaley (2008)~; Hao, Kalay et Mayhew (2010).
\end{strategyfile}

\begin{strategyfile}{Arbitrage de volatilité entre options sur indice et sur ETF}\label{strat:hf:options-cross-venue-and-volatility-arbitrage-at-speed:indexetf}
\sfield{Qui paie, et pourquoi} Les teneurs de marché de la cote dont la surface est en retard, et les utilisateurs qui paient pour les caractéristiques d'une cotation.
\sfield{Instruments et plateformes} Options sur indice et options du fonds sur le même indice.
\sfield{Signal} Écart de volatilité implicite, corrigé du style d'exercice, des dividendes et du règlement-livraison.
\sfield{Dimensionnement et exécution} Coter ou prendre la cotation en retard~; couvrir le vega sur l'autre.
\sfield{Coûts} Fourchettes des deux cotations~; erreur de modèle des corrections.
\sfield{Comment elle s'éteint} Exercice anticipé et surprises de dividende dans les options du fonds.
\sfield{Horizon, capacité, infrastructure} De quelques secondes à plusieurs jours.
\sfield{Backtester honnêtement} Les deux surfaces sur une même horloge~; les corrections telles que connues à l'époque.
\sfield{Sources} Chapitre 8~; le présent chapitre.
\end{strategyfile}

\section{Tutoriel~: les calls que personne n'a exercés}\label{tut:hf:options-cross-venue-and-volatility-arbitrage-at-speed:tut}

\begin{tutorial}
\textbf{Objectif.} Parcourir une chaîne sur plusieurs plateformes à la recherche de violations de parité et mesurer combien survivent~; pricer une box et un \omterm{def:hf:options-cross-venue-and-volatility-arbitrage-at-speed:jelly}{jelly roll}~; dimensionner le
dividend play. \textbf{État final~:} la \cref{fig:hf:options-cross-venue-and-volatility-arbitrage-at-speed:dividend} et les chiffres du texte.
\begin{tutsteps}
\item \textbf{Le scan} des conversions et des reversals sur les plateformes avec \texttt{firm.parity}.
\omcode{code/firm/optarb/firm_optarb.py}{79}{109}{Meilleurs cours sur les plateformes, frais des trois jambes et second regard après un rafraîchissement.}\label{lst:hf:options-cross-venue-and-volatility-arbitrage-at-speed:scan}
\item \textbf{Boxes et \omterm{def:hf:options-cross-venue-and-volatility-arbitrage-at-speed:jelly}{jelly rolls}}.
\omcode{code/firm/optarb/firm_optarb.py}{38}{52}{Taux implicite de la box et fourchettes~; le portage implicite de deux forwards synthétiques.}\label{lst:hf:options-cross-venue-and-volatility-arbitrage-at-speed:box}
\item \textbf{Le dividend play} pour chaque action des détenteurs défaillants (\texttt{hf\_optarb.dividend})~; le seuil d'exercice avec
\texttt{firm.american}.
\end{tutsteps}
\textbf{Pour aller plus loin.} Ajouter les frais d'emprunt aux reversals et voir quelles violations subsistent~; remplacer le bruit par des cotations qui suivent l'action avec un
délai propre à chaque plateforme~; estimer la part d'échecs à partir des positions ouvertes avant et après les dates de détachement passées.
\end{tutorial}

\section{Construction~: le module d'arbitrage sur options}\label{bld:hf:options-cross-venue-and-volatility-arbitrage-at-speed:optarb}

\begin{build}
\bfield{Rôle} Scanner les options cotées à la recherche de violations de parité, de box et de portage sur les plateformes, et dimensionner le dividend play.
\bfield{Interface} \texttt{box\_rate}, \texttt{box\_edge}, \texttt{jelly\_roll}, \texttt{Listings(n\_venues, strikes, seed)} avec \texttt{refresh}
et \texttt{best}, \texttt{scan(listings, spot\_bid, spot\_ask, fee, lag\_rho, rounds)}, \texttt{dividend\_curve(fails)},
\texttt{exercise\_threshold}. Construit sur \texttt{firm.parity} (livre~1), \texttt{firm.american} et \texttt{firm.optmm} (livre~5).
\bfield{Règles} Dollars par action~; frais par jambe~; le gain et les frais du dividend play par action, multipliés par 100 par contrat.
\bfield{Tests d'acceptation} \texttt{code/firm/optarb/tests/}~: taux et fourchettes de la box à la main~; le \omterm{def:hf:options-cross-venue-and-volatility-arbitrage-at-speed:jelly}{jelly roll} retrouve le portage~; aucune violation sans bruit~;
un bruit plus persistant laisse survivre davantage de violations~; le play ne rapporte rien en dessous de son seuil de rentabilité, et davantage à mesure que davantage de détenteurs défaillent.
\bfield{Pour aller plus loin} Emprunts et actions difficiles à emprunter~; délais des plateformes~; volatilité du SPX contre celle du SPY.
\end{build}

\begin{omsources}
\item R.~Pool, H.~R.~Stoll, R.~E.~Whaley, Failure to exercise call options: an anomaly and a trading game, \emph{Journal of Financial Markets}
11(1), 2008, 1--35.
\item J.~Hao, A.~Kalay, S.~Mayhew, Ex-dividend arbitrage in option markets, \emph{Review of Financial Studies} 23(1), 2010, 271--303.
\item J.~H.~van Binsbergen, W.~F.~Diamond, M.~Grotteria, Risk-free interest rates, \emph{Journal of Financial Economics} 143(1), 2022, 1--29.
\end{omsources}

\section{Exercices}

\begin{exercise}[$\star$]\label{exo:hf:options-cross-venue-and-volatility-arbitrage-at-speed:1}
Une box 95--105 à échéance de trois mois se traite à 9.90~\$. Quel taux implique-t-elle~?
\end{exercise}

\begin{exercise}[$\star$]\label{exo:hf:options-cross-venue-and-volatility-arbitrage-at-speed:2}
Call acheteur 4.60, put vendeur 4.10, action vendeur 100.01, strike 100, trois mois, taux 4~\%, dividendes valant 0.50 avant l'échéance. Quel est l'edge de la conversion~?
\end{exercise}

\begin{exercise}[$\star$]\label{exo:hf:options-cross-venue-and-volatility-arbitrage-at-speed:3}
Avec des frais de transaction de 0.2 cent et des frais d'exercice de 0.1 cent par action, et un gain de 30 cents par call non assigné, quelle part d'échecs
rend le dividend play rentable~?
\end{exercise}

\begin{exercise}[$\star\star$]\label{exo:hf:options-cross-venue-and-volatility-arbitrage-at-speed:4}
Pourquoi davantage de plateformes produisent-elles davantage de violations de parité, et pourquoi ne produisent-elles pas davantage de profit~?
\end{exercise}

\begin{exercise}[$\star\star$]\label{exo:hf:options-cross-venue-and-volatility-arbitrage-at-speed:5}
Pourquoi le taux implicite d'une box est-il supérieur au taux du Trésor de même maturité~?
\end{exercise}

\begin{exercise}[$\star\star$]\label{exo:hf:options-cross-venue-and-volatility-arbitrage-at-speed:6}
Que fait le dividend play au volume déclaré des options, et pourquoi cela importe-t-il aux autres~?
\end{exercise}

\begin{exercise}[$\star\star\star$]\label{exo:hf:options-cross-venue-and-volatility-arbitrage-at-speed:7}
\emph{Programmation.} Avec \texttt{firm.american.dividend\_threshold}, trouver le dividende au-delà duquel un call de strike 40 sur une action à 60~\$, à un mois de l'échéance,
doit être exercé, à une volatilité de 25~\% et un taux de 4~\%.
\end{exercise}

\begin{exercise}[$\star\star\star$]\label{exo:hf:options-cross-venue-and-volatility-arbitrage-at-speed:8}
\emph{Trouver la faille.} «~Notre scanner a trouvé 257 violations de parité pour mille instantanés sur six plateformes, un cent chacune~: un revenu régulier.~»
\end{exercise}

\section{Problème~: Les calls que personne n'a exercés}

\begin{problem}[{Problème du week-end --- les calls que personne n'a exercés}]\label{pb:hf:options-cross-venue-and-volatility-arbitrage-at-speed:1}
Un desk d'arbitrage sur options fait tourner des scans de parité, des trades de box et le dividend play.

\medskip\noindent\textbf{Partie I --- Parité.}
\begin{enumerate}
\item Définir l'\omterm{def:hf:options-cross-venue-and-volatility-arbitrage-at-speed:conversion}{arbitrage de conversion} et l'\omterm{def:hf:options-cross-venue-and-volatility-arbitrage-at-speed:conversion}{arbitrage de reversal}.
\item Décrire le scan et ses taux de violation selon le nombre de plateformes.
\item Combien survivent à un rafraîchissement, et qu'est-ce qui le décide~?
\item Pourquoi les reversals doivent-ils inclure l'emprunt~?
\end{enumerate}

\medskip\noindent\textbf{Partie II --- Taux et portage.}
\begin{enumerate}[resume]
\item Quel taux l'encadré d'exemple implique-t-il, et que rapporte son achat à 9.89~\$~?
\item Que ont trouvé van Binsbergen, Diamond et Grotteria~?
\item Définir un \omterm{def:hf:options-cross-venue-and-volatility-arbitrage-at-speed:jelly}{jelly roll} et sa valeur dans la chaîne du chapitre.
\item Qu'est-ce qui fait bouger un \omterm{def:hf:options-cross-venue-and-volatility-arbitrage-at-speed:jelly}{jelly roll} en premier~?
\end{enumerate}

\medskip\noindent\textbf{Partie III --- Le jeu du dividende.}
\begin{enumerate}[resume]
\item Quand exercer un call avant la date de détachement est-il optimal~?
\item Que ont trouvé Pool, Stoll et Whaley~?
\item Comment fonctionne le jeu, et qu'ont ajouté Hao, Kalay et Mayhew~?
\item Donner la part d'échec au seuil de rentabilité.
\end{enumerate}

\medskip\noindent\textbf{Partie IV --- Le verdict.}
\begin{enumerate}[resume]
\item Énoncer le \emph{résultat nommé}~: le profit espéré par contrat du dividend play en fonction de la part des détenteurs qui n'exercent pas.
\item Qui perd dans le dividend play, et pourraient-ils l'éviter~?
\item Que feraient des plafonds de frais plus bas au jeu~?
\item Pourquoi les trades de volatilité en cotation croisée sont-ils des trades lead--lag~?
\item Quelle fiche de stratégie porte le plus de capital~?
\item Comment estimeriez-vous la part d'échec avant une date de détachement~?
\item Que doit une société de courtage aux clients qui détiennent des calls avant une date de détachement~?
\item En une phrase~: qui paie pour les calls que personne n'a exercés~?
\end{enumerate}
\end{problem}

\section{Questions d'entretien}

\begin{interviewq}[$\star$ \iqroles{trader}]\label{iq:hf:options-cross-venue-and-volatility-arbitrage-at-speed:1}
Écrivez la parité call-put pour une action versant un dividende et dites ce que vous échangeriez si le call était trop cher d'un centime.
\end{interviewq}

\begin{interviewq}[$\star\star$ \iqroles{chercheur}]\label{iq:hf:options-cross-venue-and-volatility-arbitrage-at-speed:2}
Comment mesureriez-vous le rendement d'opportunité des bons du Trésor à partir des prix d'options~?
\end{interviewq}

\begin{interviewq}[$\star\star$ \iqroles{trader}]\label{iq:hf:options-cross-venue-and-volatility-arbitrage-at-speed:3}
Demain, une action détache un dividende de 50 cents~; vous êtes vendeur de 1\,000 calls profondément dans la monnaie. Qu'attendez-vous pendant la nuit, et que faites-vous~?
\end{interviewq}

\begin{interviewq}[$\star\star$ \iqroles{développeur}]\label{iq:hf:options-cross-venue-and-volatility-arbitrage-at-speed:4}
Concevez un scanner de parité sur 5\,000 séries réparties sur 16 plateformes, qui ne signale que les violations exécutables.
\end{interviewq}

\begin{interviewq}[$\star\star$ \iqroles{risque}]\label{iq:hf:options-cross-venue-and-volatility-arbitrage-at-speed:5}
Un desk de reversal est vendeur de 2 millions d'actions d'une action qui devient difficile à emprunter. Que se passe-t-il~?
\end{interviewq}

\begin{interviewq}[$\star\star\star$ \iqroles{chercheur}]\label{iq:hf:options-cross-venue-and-volatility-arbitrage-at-speed:6}
Dérivez la taille optimale du dividend play $q^\ast=\bigl(\sqrt{aP/c}-P\bigr)/2$ à partir de son profit espéré $aq/(P+2q)-cq$.
\end{interviewq}
